\documentclass[12pt,a4paper]{exam}
\usepackage[utf8]{inputenc}
\usepackage{amsmath,amssymb,amsfonts}
\usepackage{geometry}
\usepackage{enumitem}
\usepackage{graphicx}
\usepackage{hyperref}
\usepackage{xcolor}
\geometry{a4paper, top=2cm, bottom=2cm, left=2.5cm, right=2.5cm}
\hypersetup{colorlinks=true, linkcolor=blue, urlcolor=blue}
\renewcommand{\thequestion}{\textbf{\arabic{question}}}
\renewcommand{\questionlabel}{\thequestion.}
\pointname{ marks}
\pointsdroppedatright
\bracketedpoints
\setlength{\parindent}{0pt}
\setlength{\emergencystretch}{3em}
\allowdisplaybreaks
\newsavebox{\schmathbox}
\newcommand{\fitmath}[1]{%
  \sbox{\schmathbox}{$\displaystyle #1$}%
  \ifdim\wd\schmathbox>\linewidth
    \resizebox{\linewidth}{!}{\usebox{\schmathbox}}%
  \else
    \usebox{\schmathbox}%
  \fi}

\begin{document}

\thispagestyle{empty}
\begin{center}
  \textbf{\LARGE Diag} \\[0.3cm]
  \textbf{Time Allowed: 3 Hours} \hfill \textbf{Maximum Marks: 80} \\[0.3cm]
  \rule{\textwidth}{0.5pt}
\end{center}

\vspace{0.3cm}

\noindent\textbf{General Instructions:}
\begin{enumerate}
  \item {\normalfont\normalcolor This question paper contains 40 questions. All questions are compulsory.}
  \item {\normalfont\normalcolor Questions are grouped into sections A through E.}
  \item {\normalfont\normalcolor Use of calculators is NOT permitted.}
\end{enumerate}
\bigskip
\noindent\textbf{\large Section A: Multiple Choice \& Assertion-Reason (1 mark each)}

\begin{questions}
\question[1] If $\sin \theta = \frac{1}{2}$, then the value of $\cos \theta$ is
\begin{enumerate}[label=(\Alph*)]
  \item (A) 1
  \item (B) 0
  \item $(C) \frac{\sqrt{3}}{2}$
  \item $(D) \frac{1}{\sqrt{2}}$
\end{enumerate}
\vspace{0.15cm}

\question[1] The value of $\tan 45^\circ$ is
\begin{enumerate}[label=(\Alph*)]
  \item (A) 0
  \item (B) 1
  \item $(C) \sqrt{3}$
  \item $(D) \frac{1}{\sqrt{3}}$
\end{enumerate}
\vspace{0.15cm}

\question[1] If $\cos A = \frac{4}{5}$, then $\tan A$ is
\begin{enumerate}[label=(\Alph*)]
  \item $(A) \frac{3}{4}$
  \item $(B) \frac{5}{3}$
  \item $(C) \frac{4}{3}$
  \item $(D) \frac{3}{5}$
\end{enumerate}
\vspace{0.15cm}

\question[1] What is the value of $\sin^2 30^\circ + \cos^2 30^\circ$?
\begin{enumerate}[label=(\Alph*)]
  \item (A) 0
  \item (B) 1/2
  \item (C) 1
  \item (D) 2
\end{enumerate}
\vspace{0.15cm}

\question[1] If $3 \cot \theta = 2$, then the value of $\tan \theta$ is
\begin{enumerate}[label=(\Alph*)]
  \item (A) 2/3
  \item (B) 1/3
  \item (C) 3
  \item (D) 3/2
\end{enumerate}
\vspace{0.15cm}

\question[1] The value of $\frac{\sin 18^\circ}{\cos 72^\circ}$ is
\begin{enumerate}[label=(\Alph*)]
  \item (A) 0
  \item (B) 1
  \item (C) 2
  \item (D) -1
\end{enumerate}
\vspace{0.15cm}

\question[1] If $\sec \theta = 2$, then $\theta$ is
\begin{enumerate}[label=(\Alph*)]
  \item $(A) 30^\circ$
  \item $(B) 45^\circ$
  \item $(C) 60^\circ$
  \item $(D) 90^\circ$
\end{enumerate}
\vspace{0.15cm}

\question[1] The value of $(1 + \tan^2 \theta) \cos^2 \theta$ is
\begin{enumerate}[label=(\Alph*)]
  \item (A) 1
  \item (B) 0
  \item (C) 2
  \item (D) -1
\end{enumerate}
\vspace{0.15cm}

\question[1] If $\sin A = \cos A$, then the value of $A$ is
\begin{enumerate}[label=(\Alph*)]
  \item $(A) 30^\circ$
  \item $(B) 45^\circ$
  \item $(C) 60^\circ$
  \item $(D) 0^\circ$
\end{enumerate}
\vspace{0.15cm}

\question[1] The value of $\text{cosec}^2 \theta - 1$ is equal to
\begin{enumerate}[label=(\Alph*)]
  \item $(A) \tan^2 \theta$
  \item $(B) \sin^2 \theta$
  \item $(C) \cos^2 \theta$
  \item $(D) \cot^2 \theta$
\end{enumerate}
\vspace{0.15cm}

\question[1] If $\sin \theta = \frac{1}{2}$, then the value of $\cos \theta$ is
\begin{enumerate}[label=(\Alph*)]
  \item (A) 1
  \item (B) 0
  \item $(C) \frac{\sqrt{3}}{2}$
  \item $(D) \frac{1}{\sqrt{2}}$
\end{enumerate}
\vspace{0.15cm}

\question[1] The value of $\tan 45^\circ$ is
\begin{enumerate}[label=(\Alph*)]
  \item (A) 0
  \item (B) 1
  \item $(C) \sqrt{3}$
  \item $(D) \frac{1}{\sqrt{3}}$
\end{enumerate}
\vspace{0.15cm}

\question[1] If $\cos A = \frac{4}{5}$, then $\tan A$ is
\begin{enumerate}[label=(\Alph*)]
  \item $(A) \frac{3}{4}$
  \item $(B) \frac{5}{3}$
  \item $(C) \frac{4}{3}$
  \item $(D) \frac{3}{5}$
\end{enumerate}
\vspace{0.15cm}

\question[1] What is the value of $\sin^2 30^\circ + \cos^2 30^\circ$?
\begin{enumerate}[label=(\Alph*)]
  \item (A) 0
  \item (B) 1/2
  \item (C) 1
  \item (D) 2
\end{enumerate}
\vspace{0.15cm}

\question[1] If $3 \cot \theta = 2$, then the value of $\tan \theta$ is
\begin{enumerate}[label=(\Alph*)]
  \item (A) 2/3
  \item (B) 1/3
  \item (C) 3
  \item (D) 3/2
\end{enumerate}
\vspace{0.15cm}

\question[1] The value of $\frac{\sin 18^\circ}{\cos 72^\circ}$ is
\begin{enumerate}[label=(\Alph*)]
  \item (A) 0
  \item (B) 1
  \item (C) 2
  \item (D) -1
\end{enumerate}
\vspace{0.15cm}

\question[1] If $\sec \theta = 2$, then $\theta$ is
\begin{enumerate}[label=(\Alph*)]
  \item $(A) 30^\circ$
  \item $(B) 45^\circ$
  \item $(C) 60^\circ$
  \item $(D) 90^\circ$
\end{enumerate}
\vspace{0.15cm}

\question[1] Assertion (A): For an angle $\theta$, $\sin^2 \theta + \cos^2 \theta = 1$. Reason (R): This identity holds true for all values of $\theta$ where $0^\circ \le \theta \le 90^\circ$.
\begin{enumerate}[label=(\Alph*)]
  \item (A) Both Assertion (A) and Reason (R) are true and Reason (R) is the correct explanation of Assertion (A)
  \item (B) Both Assertion (A) and Reason (R) are true but Reason (R) is NOT the correct explanation of Assertion (A)
  \item (C) Assertion (A) is true but Reason (R) is false
  \item (D) Assertion (A) is false but Reason (R) is true
\end{enumerate}
\vspace{0.15cm}

\question[1] Assertion (A): The value of $\tan 45^\circ$ is $1$. Reason (R): $\tan \theta = \frac{\sin \theta}{\cos \theta}$ and $\sin 45^\circ = \cos 45^\circ = \frac{1}{\sqrt{2}}$.
\begin{enumerate}[label=(\Alph*)]
  \item (A) Both Assertion (A) and Reason (R) are true and Reason (R) is the correct explanation of Assertion (A)
  \item (B) Both Assertion (A) and Reason (R) are true but Reason (R) is NOT the correct explanation of Assertion (A)
  \item (C) Assertion (A) is true but Reason (R) is false
  \item (D) Assertion (A) is false but Reason (R) is true
\end{enumerate}
\vspace{0.15cm}

\question[1] Assertion (A): The value of $\sin \theta$ can be $1.5$ for some angle $\theta$. Reason (R): The range of $\sin \theta$ is $[-1, 1]$.
\begin{enumerate}[label=(\Alph*)]
  \item (A) Both Assertion (A) and Reason (R) are true and Reason (R) is the correct explanation of Assertion (A)
  \item (B) Both Assertion (A) and Reason (R) are true but Reason (R) is NOT the correct explanation of Assertion (A)
  \item (C) Assertion (A) is true but Reason (R) is false
  \item (D) Assertion (A) is false but Reason (R) is true
\end{enumerate}
\vspace{0.15cm}

\question[1] Assertion (A): $\sec^2 \theta - \tan^2 \theta = 1$. Reason (R): $1 + \tan^2 \theta = \sec^2 \theta$.
\begin{enumerate}[label=(\Alph*)]
  \item (A) Both Assertion (A) and Reason (R) are true and Reason (R) is the correct explanation of Assertion (A)
  \item (B) Both Assertion (A) and Reason (R) are true but Reason (R) is NOT the correct explanation of Assertion (A)
  \item (C) Assertion (A) is true but Reason (R) is false
  \item (D) Assertion (A) is false but Reason (R) is true
\end{enumerate}
\vspace{0.15cm}

\question[1] Assertion (A): If $\cos A = \frac{1}{2}$, then $\sin A = \frac{\sqrt{3}}{2}$. Reason (R): $\sin^2 A + \cos^2 A = 1$.
\begin{enumerate}[label=(\Alph*)]
  \item (A) Both Assertion (A) and Reason (R) are true and Reason (R) is the correct explanation of Assertion (A)
  \item (B) Both Assertion (A) and Reason (R) are true but Reason (R) is NOT the correct explanation of Assertion (A)
  \item (C) Assertion (A) is true but Reason (R) is false
  \item (D) Assertion (A) is false but Reason (R) is true
\end{enumerate}
\vspace{0.15cm}

\end{questions}
\bigskip
\noindent\textbf{\large Section B: Very Short Answer (2 marks each)}

\begin{questions}
\question[2] If $\sin(A - B) = \frac{1}{2}$ and $\cos(A + B) = \frac{1}{2}$, where $0^\circ < A + B \le 90^\circ$ and $A > B$, find $A$ and $B$.
\vspace{0.15cm}

\question[2] Show that $\tan A \cdot \tan(90^\circ - A) = 1$.
\vspace{0.15cm}

\question[2] If $3 \cot \theta = 2$, find the value of $\frac{4 \sin \theta - 3 \cos \theta}{2 \sin \theta + 6 \cos \theta}$.
\vspace{0.15cm}

\question[2] If $\sin A = \cos A$, where $0 < A < 90^\circ$, find the value of $A$.
\vspace{0.15cm}

\question[2] Evaluate $\frac{\tan 65^\circ}{\cot 25^\circ}$.
\vspace{0.15cm}

\question[2] Given $15 \cot A = 8$, find the value of $\sin A$.
\vspace{0.15cm}

\question[2] Find the value of $2 \tan^2 45^\circ + \cos^2 30^\circ - \sin^2 60^\circ$.
\vspace{0.15cm}

\question[2] Prove that $\sec^2 \theta - 1 = \tan^2 \theta$.
\vspace{0.15cm}

\end{questions}
\bigskip
\noindent\textbf{\large Section C: Short Answer (3 marks each)}

\begin{questions}
\question[3] If $\sin(A + B) = 1$ and $\cos(A - B) = \frac{\sqrt{3}}{2}$, where $0^\circ < A + B \le 90^\circ$ and $A > B$, find the values of $A$ and $B$.
\vspace{0.15cm}

\question[3] Evaluate the following expression: $\frac{2 \tan^2 45^\circ + \cos^2 30^\circ - \sin^2 60^\circ}{\sin 30^\circ + \cos 60^\circ}$
\vspace{0.15cm}

\question[3] If $7 \sin^2 \theta + 3 \cos^2 \theta = 4$, show that $\tan \theta = \frac{1}{\sqrt{3}}$.
\vspace{0.15cm}

\question[3] If $\sqrt{3} \tan \theta = 1$, find the value of $\sin^2 \theta - \cos^2 \theta$.
\vspace{0.15cm}

\question[3] Prove that $\frac{\cos A}{1 + \sin A} + \frac{1 + \sin A}{\cos A} = 2 \sec A$.
\vspace{0.15cm}

\question[3] If $\tan A = \frac{3}{4}$, find all other trigonometric ratios for angle $A$.
\vspace{0.15cm}

\end{questions}
\bigskip
\noindent\textbf{\large Section D: Long Answer (4-5 marks each)}

\begin{questions}
\question[5] Prove the following trigonometric identity: $\frac{\tan \theta}{1 - \cot \theta} + \frac{\cot \theta}{1 - \tan \theta} = 1 + \sec \theta \csc \theta$.
\vspace{0.15cm}

\question[5] If $\sec \theta + \tan \theta = p$, show that $\frac{p^2 - 1}{p^2 + 1} = \sin \theta$.
\vspace{0.15cm}

\question[5] Prove that $\frac{\tan \theta}{1 - \cot \theta} + \frac{\cot \theta}{1 - \tan \theta} = 1 + \sec \theta \csc \theta$.
\vspace{0.15cm}

\question[5] If $\sec \theta + \tan \theta = p$, prove that $\sin \theta = \frac{p^2 - 1}{p^2 + 1}$.
\vspace{0.15cm}

\end{questions}
\bigskip

\end{document}
